\documentclass[instructions.tex]{subfiles}
\begin{document}
 
\chapter{CHAPITRE XCVIII.}
 
Importance et nécessité de l'éducation chrétienne.
 
Les motifs les plus pressans engagent les parens à donner à leurs enfans une éducation vraiment chrétienne; on peut même dire que les intérêts de Dieu, les intérêts de la religion, les intérêts de la société, leurs intérêts les plus chers, et ceux de leurs enfans, reposent dans leurs mains.
 
\primo Intérêts de Dieu. Créés à l'image du Tout-Puissant, destinés à le servir sur la terre et à le posséder dans le ciel, les hommes doivent le glorifier en ce monde par une vie sainte, et en l'autre par d'éternelles louanges. Fin sublime, qui les élève au-dessus de leur propre nature, les associe aux anges eux-mêmes, et les unit étroitement à leur divin auteur.
 
Mais qui sont ceux qui, dès l'aurore de leur vie jusqu'au terme de leurs jours, tendent sans cesse vers cette fin admirable, et en remplissent constamment les devoirs les plus importans? S'il y en a, comme nous devons l'espérer, ce sont ceux-là, sans doute, qui, à l'exemple du jeune Tobie, ont été formés dès leur enfance à craindre le Seigneur, et à s'abstenir de tout péché\footnote{Quem ab infantia timere Deum docuit, et abstinere ab omni peccato. Tob. I. 10.}.
 
Imbus des vérités saintes, dès qu'ils ont pu les entendre utilement; inclinés vers le bien par de sages leçons, surtout par des exemples édifians, dès qu'ils ont été capables de les discerner; détournés du mal par des remontrances douces, des avis prudens, des corrections utiles, à mesure que les germes semblaient commencer à s'en développer en eux; enfin, préservés de la malice contagieuse du siècle, par une vigilance également soigneuse et exacte, ils ont sucé, pour ainsi dire, les principes de la vertu avec le lait de leur mère; et la grâce céleste, qui aime à se répandre sur les jeunes cœurs, secondant puissamment les efforts de leurs parens, leur enfance et leur jeunesse se sont écoulées dans la pratique de leurs devoirs; faut-il s'étonner si, dans l'âge mûr et la vieillesse, ils ne se démentent point, et ne font qu'augmenter les trésors de vertus et de mérites dont ils ont de si bonne heure posé les premiers fondemens? L'Écriture sainte nous apprend qu'il doit en être ainsi\footnote{Adolescens juxta viam suam, etiam cùm senuerit, non recedet ab eâ. Prov. 22. 6.}.
 
Or, quelle gloire ne procure point à Dieu une vie qui se passe tout entière dans l'innocence et la justice? Les anges l'admirent, les fidèles en sont édifiés, les démons en frémissent de rage; elle fait la joie des parens, le bonheur de l'ame vertueuse elle-même; et c'est un hommage continuel de soumission, de dévouement et d'amour, qu'elle rend sans cesse à la majesté divine. Aussi est-ce de cette ame généreuse que le Tout-Puissant se plaît à dire, par l'organe d'un de ses prophètes, que c'est par excellence pour sa gloire qu'il l'a créée, qu'il l'a formée, et qu'il lui a donné l'être\footnote{Omnem qui invocat nomen meum in gloriam meam creavi eum, formavi eum, et feci eum. Is. 43. 7.}.
 
Pourquoi ces exemples admirables se multiplient-ils si peu parmi nous? N'est-ce point parce que les parens, mauvais chrétiens eux-mêmes, et indolens à l'excès, au lieu d'un zèle sincère et éclairé pour l'éducation de leurs enfans, ne montrent à cet égard qu'une indifférence aussi criminelle que pernicieuse, et qu'un grand nombre s'y prennent mal et à contre sens? De là vient surtout que les enfans perdent l'inestimable trésor de l'innocence, presque aussitôt qu'ils sont en état de faire complètement usage de leur raison et de leur liberté, et qu'étant entrés de bonne heure dans la voie qui mène à la perdition, ils s'y enfoncent de plus en plus et ne s'en écartent momentanément que pour se hâter d'y rentrer. Quelle longue suite de désordres ne présage point, pour le reste de la vie, un commencement si fâcheux et si sinistre!
 
Il est vrai qu'il se rencontre dans la vie de l'homme des époques où il est vivement rappelé à son devoir; une première communion, par exemple, l'entrée dans un état fixe, une maladie grave, des graces spéciales, le ramènent quelquefois au bien. Mais si quelques-uns de ces retours sont solides, et réparent ainsi les défauts de l'éducation, les écarts de la bouillante jeunesse, combien d'autres ne sont que passagers, et laissent l'homme retomber dans sa première dépravation, en augmenter même de jour en jour les excès! Qui ne sait qu'il est des passions dont il est plus facile de se préserver entièrement, que de se guérir d'une manière parfaite et constante?
 
\secundo Intérêts de la religion. C'est dès son enfance qu'il faut commencer à inculquer à l'homme la religion. Point d'âge plus propre à y être formé. Doués de l'habitude de la foi dès leur baptême, les enfans en font aisément des actes, à mesure qu'ils apprennent à en connaître les vérités saintes. Leur esprit est docile, leur cœur maniable, leur mémoire facile; et leur imagination, semblable à une cire molle qui se prête à toutes les formes, est susceptible des plus heureuses impressions, dit saint Chrysostome. Ajoutons que ce qu'on leur dit, coule, pour ainsi parler, dans leur ame, presque sans obstacle, et que ce qu'ils ont bien appris dès cet âge, ils en conservent souvent la mémoire tout le reste de leur vie.
 
Dispositions précieuses, lesquelles étant aidées de la grace de Jésus-Christ, ne demandent qu'à être appliquées, exercées, mises soigneusement en action. Pourquoi donc si peu de parens en tirent-ils tout l'avantage qu'ils devraient en recueillir en faveur de leurs enfans? Ah! c'est que la plupart sont terrestres et mondains, ou aiment peu la religion, et la pratiquent mal, ou même la haïssent, et voudraient la voir anéantie. De là, beaucoup de pères n'en disent mot à leurs enfans, laissant tout ce soin, qui est le plus important de leurs devoirs, ou à leurs épouses, ou à des étrangers. Leur sierait-il en effet d'en parler? et pourraient-ils le faire avec avantage, tandis qu'ils n'en remplissent aucun ou presque aucun devoir?
 
O enfance chrétienne! O jeunesse de nos jours! que vous êtes à plaindre! des dangers effrayans vous menacent de toutes parts. Dans l'intérieur de vos familles, souvent nul zèle pour la gloire de Dieu, nulle pratique extérieure de religion, et beaucoup d'exemples pervers de la part de ceux qui vous ont donné le jour; au dehors, scandales multipliés, blasphèmes horribles, discours impies, livres et brochures inventés par l'enfer, compagnies séduisantes, également corruptrices de la foi et des mœurs; partout vous ne rencontrez qu'ennemis, qu'embûches et pièges. Que peuvent, hélas! contre tant d'obstacles au bien, des instituteurs sages, et des pasteurs pleins de zèle et de sollicitude?
 
Il est aisé d'expliquer, après cela, pourquoi tant de jeunes gens abandonnent les devoirs que prescrit la religion, s'enrolent même dans les rangs du vice et de l'incrédulité, presque aussitôt qu'ils sont sortis des mains de leurs instituteurs, et qu'ils ont cessé d'assister aux instructions familières de leurs pasteurs. Pères et mères insoucians, c'était surtout à vous de prévoir ces malheurs déplorables, de les prévenir à temps, et de les empêcher, à quelque prix que ce fût; votre négligence en ce point vous charge d'un compte terrible à rendre un jour au souverain Juge.
 
Mais, puisque la plupart des hommes se montrent aujourd'hui si indifférens à l'égard de la religion, et que la pratique extérieure des devoirs qu'elle prescrit semble devenir de jour en jour que le partage des personnes du sexe, à la réserve encore de quelques-unes qui déjà en secouent aussi le joug salutaire; que deviendra parmi nous cette divine institution, si c'est en vain qu'elle repose ses espérances sur les enfans qu'elle régénère et reçoit chaque jour dans son sein? Il est vrai qu'elle ne peut périr; le Tout-Puissant en a prononcé l'inébranlable stabilité jusqu'à la fin du monde; et les tempêtes affreuses auxquelles elle a survécu démontrent que le Tout-Puissant accomplit fidèlement sa promesse. Mais elle peut, et on l'a vue plus d'une fois quitter une contrée ingrate, pour se retirer dans une autre mieux disposée à la recevoir, et à profiter de ses bienfaits sans nombre.
 
Il faut cependant le dire ici; il est vrai que le zèle, la sollicitude et les soins les plus assidus peuvent quelquefois échouer malheureusement, ou qu'ils ne suffisent pas toujours pour préserver efficacement la jeunesse de tous les écueils qu'elle rencontre dans le monde, quand elle se voit forcée de s'y montrer; mais si l'on a fait tout ce qui était possible, et si l'on continue d'employer la prière, l'instance, la vigilance, tous les soins que l'on doit au plus sacré des dépôts, il faut attendre en paix le succès de sa sollicitude, et espérer que la miséricorde céleste viendra enfin arroser cette terre qui n'a point encore produit de bonnes herbes, ou qui n'en produit presque plus que de vénéneuses. Sainte Monique présente un bel exemple aux parens vraiment chrétiens. Et combien d'autres n'en pourrait-on pas citer dans le même genre!
 
\tertio Intérêts de la société. Nés pour elle, comme nos besoins, les penchans de notre cœur et le sentiment du genre humain tout entier nous le prouvent, nous devons, chacun en particulier, contribuer à son bonheur, selon nos forces, nos talens, nos facultés. Celui qui s'en montre l'ennemi, qui y trouble l'ordre, qui y répand des doctrines pernicieuses, des scandales subversifs de la religion ou des mœurs, qui y exhale le feu de la discorde, cherchant à semer la division entre les amis, les familles, les communes, les villes, etc., est un monstre, contre lequel la société a droit de s'armer du glaive de l'autorité publique. Nous ne pousserons pas plus loin ce détail.
 
Il est donc d'un grand intérêt pour la société, que ses membres soient formés, dès l'enfance, à des principes conservateurs de l'ordre, à des vertus sociales, à l'amour de Dieu, du prochain, de la patrie; et qu'on leur inspire dès leur bas-âge de l'horreur pour tout ce qui peut attirer les fléaux de la colère du Très-Haut, nuire aux autres, et leur être funeste à eux-mêmes; il est donc de l'intérêt de la société, et de son plus vif intérêt, qu'on inculque l'amour de la religion aux enfans, qu'on les en instruise de bonne heure, qu'on leur en fasse remplir les devoirs, et qu'on les détourne de tout ce qu'elle défend. La raison en est que la religion montrant à l'homme toute la noblesse de son être, toute l'étendue des devoirs qu'il a à remplir, toute la grandeur de sa destinée, pour le temps et pour l'éternité, elle présente en même temps à l'esprit, pour le conduire à son but, les lumières les plus sublimes; à son cœur, les consolations les plus douces; à ses efforts vers la vertu, les promesses les plus magnifiques; à sa faiblesse, qui est extrême, les secours les plus puissans; à ses passions bouillantes et fougueuses, le frein le plus salutaire et le plus efficace; en sorte qu'au sortir des mains de la religion, qui l'a formé, si on y ajoute la connaissance des arts et des sciences humaines, l'homme sera tout ce qu'on voudra qu'il soit, pour rendre à la société les services les plus importans, et de la manière la plus parfaite.
 
Qu'ils sont donc coupables ceux qui voudraient voir la religion bannie de la France! Si elle y était plus respectée, observée dans toute son étendue, verrait-on dans cette belle contrée les vices, les crimes, et les désordres affreux qui y règnent? qui pourrait calculer les maux de tout genre que nous a déjà attirés le discrédit dans lequel on la plonge depuis plus de soixante ans de contradictions et de combats? Retrouverait-il la France dans la France même, celui qui l'aurait bien connue il y a un siècle?
 
\quarto Intérêts des parens eux-mêmes. Fruits précieux d'une union aussi ancienne que le monde, et que le sang de Jésus-Christ a sanctifiée par un sacrement, source intarissable de graces, les enfans qu'on élève chrétiennement, et qui correspondent aux soins multipliés qu'ils reçoivent, sont comme autant de liens qui fortifient l'attachement conjugal de leurs parens; ils font couler doucement dans leurs cœurs la paix et la joie; ils prolongent leurs jours, en les rendant heureux; et ils les font vivre jusque même au-delà du trépas, parce qu'on retrouve dans la vie de ces enfans le portrait fidèle des vertus et des qualités de leurs pères et mères\footnote{Quoniam in filiis suis agnoscitur vir. Eccl. 11. 30.}.
 
Et quel genre de biens ces parens sages n'ont-ils pas lieu d'attendre des enfans dans lesquels ils voient avec délices germer et croître, sous leur main, toutes les vertus! Quelles consolations ne doivent-ils pas en espérer dans leurs peines! quels soulagemens dans leurs infirmités! quels secours dans leurs misères! quel appui dans leur vieillesse! quels soins dans leurs maladies! quelles prières ferventes et quels affectueux souvenirs après leur mort! Des enfans pieux et reconnaissans croiront-ils jamais s'être pleinement acquittés de leurs dettes envers ceux à qui, après Dieu, ils se voient redevables de tout ce qu'ils ont et de tout ce qu'ils sont?
 
Le beau spectacle que donne au monde une famille où règnent la crainte de Dieu, la piété, la subordination, la concorde, et où l'affection des parens et des enfans se confond, pour le bonheur de tous, dans une sainte et même unité\footnote{Quàm bonum et quàm jucundum habitare fratres in unum! Ps. 132. 1.}!
 
Qu'il est triste, au contraire, et qu'il est affreux, le contraste que présente une famille dont les chefs, sans piété, sans zèle, et pleins d'indolence pour l'éducation de leurs enfans, négligent de les instruire, laissent développer en eux tous les germes des passions, et loin de faire des efforts pour les former à la vertu, applaudissent bien plutôt à leurs vices, et les autorisent par leurs pernicieux exemples! Ah! si tous les genres de dépravation abondent cette horrible demeure; si les enfans s'y montrent sans crainte de Dieu, sans retenue dans leurs discours, sans frein dans leurs emportemens, sans docilité dans leur conduite; si on n'y voit qu'entêtement, que mutineries, que querelles et que guerres continuelles; si les parens n'éprouvent de la part de leurs enfans qu'insoumissions, que mépris, qu'ingratitudes, que révoltes, et même quelquefois que mauvais traitemens; si le scandale perce au dehors, et que l'infamie, la malédiction, tous les désastres viennent à la fois fondre sur ce repaire d'iniquités, de qui est-ce la faute?
 
Quelles larmes amères, quels tristes déboires se préparent pour cette vie, quel affreux désespoir et quels châtimens se ménagent pour l'éternité, les parens qui n'élèvent pas chrétiennement leurs enfans!
 
Enfin, intérêts des enfans. La vie de l'homme prend, en quelque sorte, son fondement dans l'enfance. « A cet âge, dit saint Chrysostome, elle dépend » tout entière des leçons extérieures, n'ayant encore » dans son fond presque rien qui puisse la guider; en » sorte qu'il est facile alors de l'incliner à la vertu ou » au vice. Si donc on s'y prend dès le berceau, et » comme à la venue, même des défauts puérils, pour » attirer l'homme à la vertu, l'y fortifier et lui en » faire prendre l'heureuse habitude, il ne s'en dépar- » tira pas aisément pour se jeter dans les vices con- » traires, l'habitude du bien, contractée dans l'enfance, » conservant toujours ses charmes et ses attraits\footnote{Hom. 2. in cap. 1. Joan.}. »
 
Saint Augustin décrit les efforts que font les enfans pour comprendre ce qu'on leur dit, pour saisir ce que signifient les mots qu'on emploie en leur parlant, pour en faire l'application aux objets, enfin pour les répéter et les retenir: leur éducation peut donc utilement commencer de très-bonne heure; et la raison en eux n'est point aussi tardive que beaucoup de personnes le prétendent.
 
D'ailleurs, nous apportons en naissant, outre le péché que nous avons commis dans notre premier père Adam, un levain fâcheux, dans lequel sont contenus les germes de tous les vices: germes déplorables, qui ne tardent guère à se développer; les dépits, les colères, les emportemens et les efforts de vengeance en offrent la preuve. Serait-il prudent, et conforme à l'intérêt des enfans, de laisser prendre racine, et se développer dans leurs cœurs, ces funestes germes, avant d'user de moyens prudens pour les étouffer et les détruire? Opposez-vous au principe du mal, dit un ancien; en vain chercherez-vous à l'exterminer, quand il aura pris un accroissement solide et appuyé sur la longueur du temps.
 
Les fruits d'une bonne éducation se répandent souvent sur toute la suite de la vie. Si on a le malheur de s'écarter des principes reçus dans la jeunesse, on y revient d'ordinaire tôt ou tard. Une éducation mauvaise a une influence plus grande encore, et il est difficile, il est rare même qu'on en répare tous les défauts. Combien de réprouvés s'élèveront en enfer contre leurs parens, comme ayant été les premières causes ou les fauteurs coupables de leur malheur éternel!
 
Saint Grégoire-le-Grand raconte qu'un homme, très-connu à Rome, avait un enfant qu'il aimait éperdument, et qu'il élevait fort mal. Accoutumé à blasphémer le saint nom de Dieu, dès que quelque chose s'opposait à ses désirs, cet enfant fut attaqué d'une maladie épileptique, dont il mourut. Mais peu de temps avant sa mort, étant entre les bras de son père, il fut tout-à-coup saisi de frayeur, et, tout tremblant, il se mit à crier: Défendez-moi, mon père, défendez-moi! Ensuite détournant la tête, comme pour se cacher dans le sein de son père, il répondit qu'il voyait des hommes noirs, venus pour l'enlever; et aussitôt, prononçant le blasphème dont il avait contracté l'horrible habitude, il rendit l'ame. Cet enfant était encore alors fort jeune, et il n'y avait que trois ans que le fait était arrivé quand saint Grégoire l'écrivit, sur la foi des témoins qui s'y étaient trouvés présens\footnote{Dialog. 1. 4. c. 18.}. Que d'enfans aujourd'hui blasphèment le même redoutable nom!
 
\section{Résolutions}
 
\primo Pères et mères, considérez souvent que vos enfans sont un dépôt précieux dont vous rendrez compte au souverain juge.
 
\secundo Aimez-les, vous le devez; mais pour les former et les réformer, non pour les laisser croupir dans l'ignorance, ni pour les autoriser dans leurs vices et leurs défauts; pour les sauver, non pour les perdre. 

\prayer{O mon Dieu! venez au secours de l'enfance et de la jeunesse de nos jours; et faites que le sang que Jésus-Christ a versé sur la croix pour ces tendres ames, ne soit pas perdu, et n'ait pas coulé en vain!}
 
\end{document}
